% =====================================================================
% II. SYSTEM MODEL AND DAY-AHEAD OFFERING FORMULATION
% =====================================================================
\section{System Model and Day-Ahead Offering Formulation}
\label{sec:problem_formulation}
This section models the liquid-cooled hall as a committable grid resource. We
first describe the thermal architecture and its dynamics with a virtual heat
input. We then define the cooling demand, the (offer, duration) product and
its settlement, and the safety constraints that bound the feasible
flexibility. Finally, we state the offering problem under uncertainty that
motivates the certified envelope of Section~\ref{sec:methodology}.

% ---- Fig. 1: architecture + market timeline (TikZ) ----
\begin{figure}[t]
\centering
\begin{tikzpicture}[font=\scriptsize, >=Stealth,
  box/.style={draw, rounded corners=1pt, align=center, inner sep=2.5pt},
  lbl/.style={align=center, inner sep=1pt}]
  % thermal chain (bottom row)
  \node[box, fill=red!8]    (gpu)  {GPU racks\\ $\Tj$,\ $\bar Q{+}w_t$};
  \node[box, fill=blue!8, right=6mm of gpu]  (loop) {coolant loop\\ $\Tw$};
  \node[box, fill=cyan!8, right=6mm of loop] (fac)  {facility water\\ $\Tf$};
  \node[box, fill=gray!12, right=6mm of fac] (chil) {chiller\\ $\Pcool_t$};
  \draw[->] (gpu) -- node[above, lbl]{$h_{\mathrm{jw}}$} (loop);
  \draw[->] (loop) -- node[above, lbl]{$\uext_t$} (fac);
  \draw[->] (fac) -- node[above, lbl]{$\urej_t$} (chil);
  % constraint labels (below chain)
  \node[lbl, below=1.5mm of gpu, text=red!60!black]  {$\Tj \le \Tmax$};
  \node[lbl, below=1.5mm of loop, text=blue!60!black]
    {supply $\ge \Tdew{+}\delta_{\mathrm{c}}$};
  \node[lbl, below=1.5mm of fac]  {$\Tf \le \Tf^{\max}$};
  % top row: context + market
  \node[box, fill=yellow!12, text width=21mm, above=7mm of gpu.north west,
        anchor=south west] (ctx)
    {context $c$:\\ NWP $\Tdew$ forecast,\\ job-aware load forecast};
  \node[box, fill=green!8, text width=39mm, above=7mm of chil.north east,
        anchor=south east] (mkt)
    {day-ahead market:\\ offer $(q_h,d)$ $\to$ activation $r_h$\\
     $\to$ whole-hour settlement};
  \draw[->, dashed] (ctx.east) -- (mkt.west);
  \draw[->, dashed] (chil.north) -- node[right, lbl, pos=0.45]
    {metered $\Pcool$ vs.\ $\Pbase$} (mkt.south -| chil.north);
\end{tikzpicture}
\caption{Grid-interactive liquid-cooled hall as a committable resource. The
virtual inputs $(\uext_t,\urej_t)$ move heat along the chain; the operator
commits an (offer, duration) pair day-ahead from the context-dependent
envelope and is settled on metered whole-hour cooling energy.}
\label{fig:phys_arch}
\vspace{-3mm}
\end{figure}

\vspace{-4mm}
\subsection{Liquid-Cooled Hall with a Virtual Heat Input}
\label{subsec:arch}

We consider the thermal path of a hall with $1$\,MW of IT load served by
direct-to-chip liquid cooling, as depicted in Fig.~\ref{fig:phys_arch}: GPU
cold plates reject heat into a coolant loop, a coolant distribution unit
passes it to a facility-water loop, and a chiller plant rejects it to ambient.
Thermally, the hall must keep the silicon-adjacent hotspot and the loop
temperatures within limits; electrically, the chiller plant is a
controllable load whose temporary reduction is the product being sold.

The thermal state at control step $t$ is
\begin{equation}
x_t = [\Tj[t],\; \Tw[t],\; \Tf[t]]^\top,
\end{equation}
collecting the hotspot proxy, the coolant-loop temperature, and the
facility-water temperature. The control vector is the pair of \emph{virtual
heat inputs}
\begin{equation}
u_t = [\uext_t,\; \urej_t]^\top,
\end{equation}
where $\uext_t$ is the heat-extraction rate from the coolant loop into the
facility loop and $\urej_t$ the heat-rejection rate of the chiller plant. The
virtual inputs absorb the operator's physical handles (supply-temperature
setpoint, flow) into heat-rate variables. This choice renders the dynamics
linear and confines all nonlinearity to a state-dependent input set defined
below.

\begin{remark}[Realization layer]\label{rem:realization}
A commanded extraction rate $\uext$ is realized by the supply-temperature
setpoint $T^{\mathrm{in}} = \Tw - \uext/(\mdot c_p)$ at the operating flow,
and a rejection rate $\urej$ by the chiller load setpoint. The set
\eqref{eq:Ux} below is the image of the physical actuator ranges under this
map,
so every certified virtual command corresponds to an admissible setpoint
pair; the actuators' bilinear flow--temperature coupling never enters the
commitment problem (cf.~\cite{shen2026cdro}).
\end{remark}

The disturbance over an hour is
\begin{equation}
\xi = \bigl[w_{0:N-1},\; \delta^{\mathrm{dew}}\bigr] ,
\end{equation}
where $w_t$ is the step residual of the IT heat load around its day-ahead
forecast $\bar Q$, driven by bursty, non-Gaussian AI workload dynamics, and
$\delta^{\mathrm{dew}}$ is the hour's dew-point forecast error, which shifts
the anti-condensation boundary. The day-ahead context vector
\begin{equation}
c_h = \bigl[\Tdew^{\mathrm{fc}},\ \text{workload-forecast features}\bigr]
\label{eq:context}
\end{equation}
collects what is knowable at the offering deadline; every object downstream
(envelope, disturbance sets, offer) is a function of $c_h$.

\vspace{-3mm}
\subsection{Thermal Dynamics}
\label{subsec:dynamics}

The hall is modeled by lumped energy balances on the three nodes.

\emph{Hotspot node.} The silicon-adjacent metal has small capacitance
$C_\mathrm{j}$ (time constant tens of seconds), responding to workload bursts
much faster than the loop can react:
\begin{equation}
C_\mathrm{j}\,\frac{d\Tj}{dt} = \bar Q + w(t)
 - h_{\mathrm{jw}}\bigl(\Tj(t)-\Tw(t)\bigr),
\label{eq:chip}
\end{equation}
where $h_{\mathrm{jw}}$ is the junction-to-coolant heat-transfer coefficient.

\emph{Coolant-loop node.} The loop's fluid inventory provides the first
tranche of thermal mass (time constant of minutes):
\begin{equation}
C_\mathrm{w}\,\frac{d\Tw}{dt} = h_{\mathrm{jw}}\bigl(\Tj(t)-\Tw(t)\bigr)
 - \uext(t).
\label{eq:loop}
\end{equation}

\emph{Facility-water node.} The facility loop provides the second, larger
tranche (time constant near an hour), and makes a half-hour product
physically possible:
\begin{equation}
C_\mathrm{f}\,\frac{d\Tf}{dt} = \uext(t) - \urej(t).
\label{eq:facility}
\end{equation}

The state-dependent input set is
\begin{equation}
\mathcal{U}(x;c)=\Bigl\{u:
\begin{aligned}
&\ell(\Tw) \le \uext \le \mdot c_p\bigl(\Tw-\vartheta(c)\bigr),\\
&\uext \le \mdot[f] c_p\bigl(\Tw-\Tf-\delta_{\mathrm{hx}}\bigr),\
0 \le \urej \le \bar u^{\mathrm{rej}}
\end{aligned}
\Bigr\},
\label{eq:Ux}
\end{equation}
where $\mdot[f]$ is the facility-loop flow and
$\ell(\Tw)=[\mdot_{\min} c_p(\Tw - T^{\mathrm{in}}_{\max})]_+$ is the forced
extraction once the loop exceeds the supply-temperature ceiling; the set is
\emph{affine} in $x$ for a fixed context. The upper extraction bound carries
the model's signature constraint: the supply-temperature floor
\begin{equation}
\vartheta(c) = \max\bigl\{\vartheta_{\min},\ \Tdew^{\mathrm{fc}}
 + \delta_{\mathrm{c}}\bigr\}
\label{eq:floor}
\end{equation}
sits above the dew point to preclude condensation on server boards, so
humidity directly truncates the admissible heat extraction. The second bound
is the passive heat-exchanger coupling between the loops, and the third the
chiller equipment limit. We assume the hall-air dew point tracks the ambient
dew point (outdoor-air-coupled halls without active dehumidification), with
the margin $\delta_{\mathrm{c}}$ absorbing the residual gap; halls with
tightly conditioned air would replace $\Tdew^{\mathrm{fc}}$ by an indoor
humidity forecast without changing the structure. Sub-minute chip-level bursts are summarized by
intra-interval peak statistics within $w_t$; hardware frequency scaling
remains the ultimate backstop beyond all modeled uncertainty sets and appears
nowhere else in the model.

For certification, \eqref{eq:chip}--\eqref{eq:facility} are discretized
exactly (zero-order hold) at the control interval $\Delta t=5$\,min, giving
$x_{t+1}=Ax_t+Bu_t+E(\bar Q + w_t)$; closed-loop simulation integrates the
same model at a $30$\,s step. Parameters (Table~\ref{tab:params}) are
calibrated so the plant reproduces the measured supply-temperature
flexibility of a production liquid-cooled hall ($17\to 25^\circ$C shedding
$63\%$ of cooling power~\cite{gheni2026liquid}), with loop, facility, and
junction time constants of $5$--$8$\,min, ${\sim}50$\,min, and ${\sim}24$\,s.

% ---- Table: plant parameters ----
\begin{table}[t]
\centering
\caption{Plant parameters of the calibrated $1$\,MW$_{\mathrm{IT}}$ hall.}
\label{tab:params}
\setlength{\tabcolsep}{3.5pt}
\scriptsize
\begin{tabular}{llll}
\toprule
\textbf{Symbol} & \textbf{Value} & \textbf{Symbol} & \textbf{Value} \\
\midrule
$C_\mathrm{j}$ & $1.0$\,MJ/K & $\Tmax$ & $85^\circ$C \\
$C_\mathrm{w}$ & $25$\,MJ/K & $\Tf^{\max}$ & $35^\circ$C \\
$C_\mathrm{f}$ & $60$\,MJ/K & $\delta_{\mathrm{c}}$ & $2$\,K \\
$h_{\mathrm{jw}}$ & $40$\,kW/K & $\vartheta_{\min}$ & $18^\circ$C \\
$\mdot c_p$ & $54.6$\,kW/K & $\bar u^{\mathrm{rej}}$ & $1.3$\,MW \\
$\delta_{\mathrm{hx}}$ & $2$\,K & $\uext$ ramp & $200$\,kW/min \\
$\eta$ (ref.\ COP) & $3.8$ & $T_{\mathrm{thr}}$ & $70^\circ$C \\
$\Delta t$ / $\dH$ & $5$\,min / $1$\,h & $\mdot_{\min}/\mdot$ & $0.3$ \\
\bottomrule
\end{tabular}
\vspace{-3mm}
\end{table}

\vspace{-3mm}
\subsection{Cooling Demand, Product, and Settlement}
\label{subsec:product}

Under a frozen-efficiency surrogate, the metered cooling demand is affine in
the rejection rate,
\begin{equation}
\Pcool_t = P_0 + \urej_t/\eta,
\label{eq:coolpower}
\end{equation}
where $\eta$ is the chiller coefficient of performance at the reference point
and $P_0$ collects pumping. The surrogate freezes the weather dependence of
plant efficiency within the hour; its drift enters the energy and degradation
terms, not the certificate. The no-participation trajectory defines the
exogenous baseline $\Pbase$.

For each hour $h$ the operator offers a pair $(q_h, d)$: \emph{if activated,
reduce cooling power by up to $q_h$ below baseline, sustainable for up to $d$
minutes within the hour}. The duration $d\in\{15,30,60\}$\,min is a product
parameter fixed per market (we evaluate $d=30$); activation is an exogenous
Bernoulli draw with probability $p_{\mathrm{act}}$ per committed hour, and an
activated hour exercises the full window, $r_h = d/60$ (idle hours have
$r_h=0$). The market interval is $\dH = 1$\,h.

\begin{remark}[Market realism]\label{rem:market}
The product is a \emph{generic short-duration reserve abstraction} and the
only institutional element in the model. Qualification, telemetry, and performance rules of
the ERCOT Contingency Reserve Service (ECRS) are not modeled; ECRS clearing
prices serve as the capacity-price stream, and $p_{\mathrm{act}}$ and $\gamma$ are product design parameters
whose sensitivity Section~\ref{subsec:cert} reports. Economic results are
stated \emph{under this ERCOT-price-driven stylized settlement}, not as
realized revenues of an enrolled participant.
\end{remark} Hour $h$ settles on metered energy
against the baseline:
\begin{equation}
\begin{aligned}
\text{revenue} &= \picap_h q_h \dH,\qquad
\text{penalty} = \gamma\,\picap_h\, s_h,\quad \gamma \ge 1,\\
s_h &= \Bigl[r_h q_h \dH -
\textstyle\sum_{t\in h}(\Pbase-\Pcool_t)\Delta t\Bigr]_+
\!\wedge\, r_h q_h \dH,
\end{aligned}
\label{eq:settlement}
\end{equation}
where the $\wedge$ term caps the shortfall at the obligation; the hour
additionally settles real-time energy deviations at the spot price and a
degradation cost $c_{\mathrm{deg}}\!\int[\Tj(t)-T_{\mathrm{thr}}]_+dt$ that
prices sustained hotspot elevation above the wear threshold
$T_{\mathrm{thr}}=70^\circ$C. Whole-hour metering in \eqref{eq:settlement} drives the
certification: recovery \emph{within} the activated hour cancels delivered
energy, so certifying the activation window alone over-promises; the envelope
of Section~\ref{sec:methodology} certifies both the window depth and the hour
energy.

\begin{remark}[Baseline]\label{rem:baseline}
The baseline is exogenous and known. Baseline manipulation is a documented
phenomenon in demand-response programs~\cite{chao2011demand}; it is a
separate problem from deliverability, and the baseline is taken as given
here as in the buildings-reserves literature~\cite{vrettos2016robust}.
\end{remark}

\vspace{-3mm}
\subsection{Operational Constraints and the Flexibility Envelope}
\label{subsec:constraints}

The plant must satisfy the following at every step.

\emph{C1. Hotspot safety:} $\Tj[t] \le \Tmax$, with $\Tmax=85^\circ$C the
lumped proxy below silicon throttling.

\emph{C2. Anti-condensation (hard, safety-critical):} the floor
\eqref{eq:floor} inside $\mathcal{U}(x;c)$; an excursion risks condensate on
server boards.

\emph{C3. Facility limit:} $\Tf[t] \le \Tf^{\max}$, protecting downstream
equipment.

\emph{C4. Actuation:} $u_t \in \mathcal{U}(x_t;c)$ with ramp limits on
$\uext$.

A commitment $(q,d)$ is \emph{deliverable from state $x_0$ in context $c$} if
some admissible input sequence satisfies C1--C4 together with the two product
clauses: window depth $\Pcool_t \le \Pbase - q$ during the activation window,
and whole-hour energy per \eqref{eq:settlement}. The \emph{flexibility
envelope}
\begin{equation}
F(x,c) = \max\bigl\{q : (q,d)\ \text{deliverable from}\ x\ \text{in}\ c\bigr\}
\label{eq:envelope_def}
\end{equation}
is the maximal committable depth. It shrinks with humidity through
\eqref{eq:floor}, with workload volatility through the buffer that bursts
consume, and with duration $d$; a fixed-margin commitment is therefore either
unsafe or wasteful across contexts. Context-dependent robustness resolves the
same dilemma in the companion control problem~\cite{shen2026cdro}; here it
enters at the market layer.

\vspace{-4mm}
\subsection{Day-Ahead Offering Problem Under Uncertainty}
\label{subsec:problem}

At the day-ahead deadline the operator chooses commitments
$\{q_h\}_{h=1}^{24}$ knowing only the context \eqref{eq:context}; on the
operating day, activations arrive and the realized disturbances
$\xi_t$ act on the plant under a causal real-time policy $\pi$. The ideal
stochastic offering problem is
\begin{equation}
\begin{aligned}
\max_{\{q_h\},\,\pi}\ \
&\mathbb{E}_{\xi\sim\mathbb{P}^*}\Bigl[\textstyle\sum_h
 \picap_h q_h \dH - \gamma\,\picap_h s_h - C^{\mathrm{rt}}_h - D_h \Bigr]\\
\text{s.t.}\ \ & x_{t+1}=Ax_t+Bu_t+E(\bar Q+w_t),\quad u_t=\pi_t(\cdot),\\
& \text{C1--C4 for all } t,
\end{aligned}
\label{eq:stochastic_opt}
\end{equation}
where $C^{\mathrm{rt}}_h$ and $D_h$ are the real-time energy and degradation
costs and $\mathbb{P}^*$ is the unknown joint law of workload and dew-point
residuals. Three features make the problem hard: $\mathbb{P}^*$ is unknown,
nonstationary, and context-dependent; safety (C1--C3) admits no monetary
compensation while shortfalls are merely priced, so the two consequences
carry different stakes; and the commitment precedes delivery by hours, so
all robustness must be decided from the context alone.
Section~\ref{sec:methodology} resolves this with a certified envelope
robustified at two nested confidence levels.

\vspace{-2mm}
